% PRINCIPIA PHYSICA -- pilot topic 2
% Hamiltonian reconstruction as a physical theory object
% arXiv-style preprint, self-contained (no external .bib required)
\documentclass[10pt,a4paper]{article}

\usepackage[margin=0.9in]{geometry}
\usepackage{amsmath,amssymb,amsthm,mathtools}
\usepackage{newtxtext}
\usepackage[varvw]{newtxmath}
\usepackage{bm}
\usepackage{enumitem}
\usepackage{listings}
\usepackage{titlesec}
\usepackage[hidelinks]{hyperref}

\titlespacing*{\section}{0pt}{1.4ex plus .2ex}{0.7ex}
\titlespacing*{\subsection}{0pt}{1.1ex plus .2ex}{0.4ex}
\titleformat{\section}{\normalfont\large\bfseries}{\thesection}{0.6em}{}
\titleformat{\subsection}{\normalfont\bfseries}{\thesubsection}{0.6em}{}

% Compact display spacing, reapplied at every size change.
\makeatletter
\g@addto@macro\normalsize{%
  \setlength{\abovedisplayskip}{5pt plus 2pt minus 2pt}%
  \setlength{\belowdisplayskip}{5pt plus 2pt minus 2pt}%
  \setlength{\abovedisplayshortskip}{2pt plus 1pt}%
  \setlength{\belowdisplayshortskip}{3pt plus 1pt}}
\makeatother

% Tighter vertical spacing around the many typed statement environments.
\makeatletter
\def\thm@space@setup{\thm@preskip=4pt plus1pt minus1pt
                     \thm@postskip=4pt plus1pt minus1pt}
\makeatother

\lstset{
  basicstyle=\ttfamily\scriptsize,
  columns=fullflexible,
  keepspaces=true,
  frame=single,
  framesep=4pt,
  xleftmargin=6pt,
  breaklines=true
}

% ---------------------------------------------------------------
% Epistemic environments. Every numbered statement in this paper
% carries exactly one epistemic type. The types share a counter so
% that the type of a claim can never be inferred from its number.
% ---------------------------------------------------------------
\theoremstyle{definition}
\newtheorem{definition}{Definition}[section]
\newtheorem{axiom}[definition]{Axiom}
\newtheorem{assumption}[definition]{Assumption}
\newtheorem{conjecture}[definition]{Conjecture}
\newtheorem{interpretation}[definition]{Interpretation}
\newtheorem{empirical}[definition]{Empirical statement}
\newtheorem{established}[definition]{Established result}
\newtheorem{newresult}[definition]{New result claimed by this work}
\newtheorem{example}[definition]{Example}
\newtheorem{nonexample}[definition]{Non-example}
\newtheorem{remark}[definition]{Remark}

\theoremstyle{plain}
\newtheorem{lemma}[definition]{Lemma}
\newtheorem{proposition}[definition]{Proposition}
\newtheorem{theorem}[definition]{Theorem}
\newtheorem{corollary}[definition]{Corollary}

\newcommand{\Rr}{\mathbb{R}}
\newcommand{\Cinf}[1]{C^{\infty}(#1)}
\newcommand{\Xf}{\mathfrak{X}}
\newcommand{\Ham}{\mathrm{Ham}}
\newcommand{\Symp}{\mathrm{Symp}}
\newcommand{\dR}{\mathrm{dR}}
\newcommand{\Th}{\mathbb{T}}
\DeclareMathOperator{\Ker}{ker}
\DeclareMathOperator{\Imm}{im}
\DeclareMathOperator{\dive}{div}
\DeclareMathOperator{\Area}{Area}

\title{\bfseries Hamiltonian Reconstruction as a Physical Theory Object\\[2mm]
\large A typed reconstruction of finite-dimensional Hamiltonian mechanics,\\
with its composition laws, its empirical interpretation, and its limits}

\author{PRINCIPIA PHYSICA pilot, Topic 2\\
\small Compositional Foundations of Finite-Dimensional Classical Mechanics}

\date{16 August 2026}

\begin{document}
\sloppy
\hbadness=10000
\maketitle

\begin{abstract}
We reconstruct finite-dimensional Hamiltonian mechanics as an explicitly
\emph{typed} object --- a kinematic type, an observable type, a generator, a
dynamical law, and an empirical interpretation --- under six stated axioms and
one stated assumption. After proving the standard structural facts within this
presentation, we use the typing to isolate four results bearing on the
compositional thesis under test in this pilot. (i) The Cartesian monoidal
product of Hamiltonian theory objects has image exactly the non-interacting
composites, so interaction is provably outside the range of $\otimes$. (ii)
Symplectic morphisms are immersions, so coarse-graining is not a morphism of
the category and reduction must be a different operation. (iii) A globally
asymptotically stable dissipative field on a surface is not Hamiltonian for
\emph{any} symplectic form, not merely for the standard one. (iv) A matched
pair of counterexamples separates mathematical from empirical equivalence in
both directions, by an explicit Gr\"onwall bound in one direction and an
explicit period-function obstruction in the other. We state four structural
limits --- Darboux triviality, the failure of the symplectic ``category'' to
be a category, the Groenewold--van Hove obstruction, and symplectic rigidity
--- and give falsification conditions. We do not claim the compositional
thesis is established, only that this reconstruction makes it precise enough
to be attacked. A Haskell appendix renders the type skeleton; it is a typing,
not a formal proof.
\end{abstract}

% ===============================================================
\section{Introduction and epistemic conventions}
% ===============================================================

\subsection{The thesis under test}

The PRINCIPIA PHYSICA pilot investigates the thesis that \emph{physics is the
study of empirically realized compositional mathematical structures}. This
paper does not assume that thesis. It performs one bounded task on which the
thesis can be tested: it reconstructs finite-dimensional Hamiltonian mechanics
as a typed object, records exactly which of its features are mathematical and
which are empirical, and reports where the compositional reading succeeds,
where it is silent, and where it provably fails.

Three separations organise the paper.
\textbf{(S1) Structure versus law:} the kinematic type $(M,\omega)$ and the
dynamical law $\iota_{X_H}\omega = dH$ are logically independent data, and
Section~\ref{sec:reconstruction} shows what each contributes.
\textbf{(S2) Mathematical versus physical equivalence:} an isomorphism of the
underlying mathematical objects is neither necessary nor sufficient for
empirical indistinguishability, and Section~\ref{sec:empirical} proves both
non-implications.
\textbf{(S3) Composition versus interaction:} the monoidal product of
Hamiltonian systems exists, and Section~\ref{sec:composition} proves that it
cannot express interaction, so any compositional account of physics must
supply an interconnection operation that is not $\otimes$.

\subsection{Scope}

Everything below is confined to smooth, finite-dimensional, real classical
systems. Field theory, infinite-dimensional analysis, general relativity,
quantum theory, and quantum field theory appear only as \emph{boundary tests}:
they are named to mark where the present type stops applying, and no result is
claimed inside them. No new experimental claim is made anywhere in this paper.

\subsection{Epistemic conventions}
\label{sec:conventions}

Every numbered statement carries exactly one of the following types, and the
types share a single counter so that a statement's number never encodes its
epistemic status:
\emph{Definition}, \emph{Axiom}, \emph{Assumption}, \emph{Conjecture},
\emph{Lemma}, \emph{Proposition}, \emph{Theorem}, \emph{Corollary},
\emph{Interpretation}, \emph{Empirical statement}, \emph{Established result},
\emph{New result claimed by this work}. We use them as follows.
\emph{Axiom} marks a stipulation constitutive of the object being defined; it
is not asserted to be true of the world. \emph{Assumption} marks a restriction
adopted to make a proof go through; assumptions are never upgraded to theorems
anywhere in this paper, and where one fails we give a non-example.
\emph{Established result} marks a statement proved elsewhere and cited, not
reproved here. \emph{New result claimed by this work} marks statements we
prove here and put forward as the contribution; where a statement is
elementary and probably folklore we say so in the same breath, and the claim
is then to the \emph{explicit statement} rather than to mathematical novelty.
\emph{Empirical statement} marks a claim about measurement outcomes, either
cited to a measurement or flagged as an input. \emph{Interpretation} marks a
semantic bridge between formal structure and physical vocabulary;
interpretations are not derived. All references were checked for bibliographic
identity before citation, and Appendix~\ref{app:verification} records what was
confirmed for each and what remains inferred.

% ===============================================================
\section{The typed object}
\label{sec:object}
% ===============================================================

\subsection{Types}

\begin{definition}[Type vocabulary]
\label{def:types}
We use five sorts. A \emph{kinematic type} is a pair $(M,\omega)$ with $M$ a
smooth manifold and $\omega$ a differential $2$-form on $M$. An
\emph{observable type} over $(M,\omega)$ is the commutative $\Rr$-algebra
$\mathcal{A}(M) := \Cinf{M,\Rr}$. A \emph{generator} is an element
$H \in \mathcal{A}(M)$. A \emph{process type} over $(M,\omega)$ is a smooth
local flow on $M$. An \emph{interpretation} is the data specified in
Definition~\ref{def:interpretation}.
\end{definition}

\begin{definition}[Hamiltonian theory object]
\label{def:object}
A \emph{Hamiltonian theory object} is a tuple
\[
  \Th \;=\; \bigl(M,\ \omega,\ \mathcal{A}(M),\ H,\ \Phi,\ \mathcal{I}\bigr)
\]
satisfying Axioms \ref{ax:kin}--\ref{ax:interp} below. We write
$\Th^{\mathrm{math}} = (M,\omega,\mathcal{A}(M),H,\Phi)$ for the
\emph{bare} object obtained by discarding the interpretation, and
$U: \Th \mapsto \Th^{\mathrm{math}}$ for that forgetting operation.
\end{definition}

\subsection{Axioms}

\begin{axiom}[Kinematic type]
\label{ax:kin}
$M$ is a smooth, second-countable, Hausdorff manifold of finite dimension
$\dim M = 2n$ with $n \ge 0$. Every substantive statement below has $n\ge 1$;
the case $n=0$ is admitted solely so that the one-point object can serve as
the monoidal unit in Corollary~\ref{cor:secondop}.
\end{axiom}

\begin{axiom}[Symplectic structure]
\label{ax:symp}
$\omega \in \Omega^2(M)$ satisfies $d\omega = 0$ and is nondegenerate: for
each $x \in M$ the map
$\omega^{\flat}_x : T_xM \to T^{*}_xM$, $v \mapsto \omega_x(v,\cdot\,)$,
is a linear isomorphism.
\end{axiom}

\begin{axiom}[Observable type]
\label{ax:obs}
The observables are $\mathcal{A}(M) = \Cinf{M,\Rr}$ with pointwise
operations. Physical quantities are elements of $\mathcal{A}(M)$; the value of
the quantity $f$ in the state $x$ is the real number $f(x)$.
\end{axiom}

\begin{axiom}[Generator]
\label{ax:gen}
A distinguished element $H \in \mathcal{A}(M)$ is given.
\end{axiom}

\begin{axiom}[Dynamical law]
\label{ax:law}
The evolution $\Phi$ is the maximal local flow of the unique vector field
$X_H$ determined by $\iota_{X_H}\omega = dH$, and evolution is autonomous:
$\Phi_{t+s} = \Phi_t \circ \Phi_s$ wherever both sides are defined.
\end{axiom}

\begin{axiom}[Interpretation]
\label{ax:interp}
An interpretation $\mathcal{I}$ in the sense of
Definition~\ref{def:interpretation} is given.
\end{axiom}

Axioms~\ref{ax:kin}--\ref{ax:interp} are stipulations. Nothing in this paper
asserts that any physical system satisfies them; Section~\ref{sec:empirical}
addresses when the assertion is warranted and Section~\ref{sec:nonexamples}
exhibits systems for which each of Axioms~\ref{ax:kin}, \ref{ax:symp} and
\ref{ax:law} fails.

\begin{assumption}[Completeness]
\label{as:complete}
Where a statement quantifies over all $t \in \Rr$ we assume $X_H$ is complete,
so that $\Phi$ is a global $\Rr$-action. This assumption is used only where
stated.
\end{assumption}

Assumption~\ref{as:complete} is a genuine restriction, not a technicality.

\begin{nonexample}[Incompleteness]
\label{nex:incomplete}
On $(\Rr^2,\ \omega_0 = dq\wedge dp)$ take $H(q,p)=\tfrac12 p^2 - q^4$. On the
level set $H = E$ one has $\dot q = p = \pm\sqrt{2(E+q^4)}$ along the branch
$p>0$, so the time to reach $q = +\infty$ from $q_0$ with $E + q_0^4 > 0$ is
\[
  \int_{q_0}^{\infty}\frac{dq}{\sqrt{2\,(E+q^4)}} \;<\; \infty ,
\]
because the integrand is $O(q^{-2})$ as $q \to \infty$. Hence $X_H$ is
incomplete: the solution leaves every compact set in finite time. This object
satisfies Axioms~\ref{ax:kin}--\ref{ax:law} and violates
Assumption~\ref{as:complete}.
\end{nonexample}

\subsection{The interpretation}

\begin{definition}[Interpretation, empirical model]
\label{def:interpretation}
An \emph{interpretation} of a bare object $\Th^{\mathrm{math}}$ is a tuple
$\mathcal{I} = (P,\ \mu,\ \mathcal{U},\ \varepsilon,\ T)$ where
\begin{itemize}[leftmargin=*,itemsep=1pt]
\item $P \subseteq M$ is the set of \emph{preparable} states, the image of a
      laboratory preparation procedure;
\item $\mu : M \to \Rr^{k}$ is a smooth \emph{readout map}, the composite of
      the quantities actually recorded by the apparatus;
\item $\mathcal{U}$ assigns a physical dimension and unit to each component of
      $\mu$ and to $H$;
\item $\varepsilon > 0$ is the readout resolution and $T > 0$ the accessible
      time window.
\end{itemize}
The \emph{empirical model} of $\Th$ is
\[
  \mathrm{Emp}(\Th)\;=\;
  \bigl\{\,(z,t,\mu(\Phi_t z)) \;:\; z \in P,\ t\in[0,T] \,\bigr\}
  \ \subseteq\ P\times[0,T]\times\Rr^{k},
\]
recorded together with $\mathcal{U}$ and read to precision $\varepsilon$.
\end{definition}

\begin{interpretation}[Locus of empirical content]
\label{int:locus}
The empirical content of $\Th$ resides in $\mathrm{Emp}(\Th)$, not in
$\Th^{\mathrm{math}}$. Two consequences are used below: (i) any feature of
$(M,\omega,H)$ invisible to $\mu$ on $P$ within $\varepsilon$ over $[0,T]$ is,
for that interpretation, empirically idle; (ii) $\mathcal{I}$ is not
recoverable from $\Th^{\mathrm{math}}$, because $U$ discards it, which is what
Proposition~\ref{prop:isonotphys} makes precise.
\end{interpretation}

% ===============================================================
\section{Reconstruction: what the axioms generate}
\label{sec:reconstruction}
% ===============================================================

\subsection{The dynamical law is well posed}

\begin{lemma}[Existence and uniqueness of $X_H$]
\label{lem:XH}
Under Axioms~\ref{ax:kin}--\ref{ax:gen} there is exactly one smooth vector
field $X_H$ on $M$ with $\iota_{X_H}\omega = dH$.
\end{lemma}

\begin{proof}
By Axiom~\ref{ax:symp} the bundle map
$\omega^{\flat}:TM \to T^{*}M$ is a fibrewise linear isomorphism, and it is
smooth because $\omega$ is. A smooth bundle isomorphism has a smooth inverse
$(\omega^{\flat})^{-1}$. Setting $X_H := (\omega^{\flat})^{-1}(dH)$ gives a
smooth vector field with $\iota_{X_H}\omega = dH$. If $X$ and $X'$ both
satisfy the equation then $\iota_{X-X'}\omega = 0$, and nondegeneracy forces
$X = X'$ pointwise.
\end{proof}

\begin{definition}[Poisson bracket]
\label{def:bracket}
For $f,g \in \mathcal{A}(M)$ set $\{f,g\} := \omega(X_f,X_g)$.
\end{definition}

\begin{lemma}[Bracket identities]
\label{lem:poisson}
For all $f,g,h \in \mathcal{A}(M)$:
\begin{enumerate}[label=(\roman*),leftmargin=*,itemsep=1pt]
\item $\{f,g\} = X_g(f) = -\{g,f\}$;
\item $\{fg,h\} = f\{g,h\} + g\{f,h\}$ (Leibniz rule);
\item $\iota_{[X,Y]}\omega = d\bigl(\omega(Y,X)\bigr)$ for all
      $X,Y$ with $\mathcal{L}_X\omega = \mathcal{L}_Y\omega = 0$;
\item $[X_f,X_g] = -X_{\{f,g\}}$;
\item $\{f,\{g,h\}\}+\{g,\{h,f\}\}+\{h,\{f,g\}\} = 0$.
\end{enumerate}
Hence $(\mathcal{A}(M),\cdot,\{\cdot,\cdot\})$ is a Poisson algebra.
\end{lemma}

\begin{proof}
(i) $\omega(X_f,X_g) = (\iota_{X_f}\omega)(X_g) = df(X_g) = X_g(f)$, and
antisymmetry is that of $\omega$.
(ii) $X_h$ is a derivation of $\mathcal{A}(M)$ and $\{f,h\} = X_h(f)$ by (i).
(iii) Cartan's formula gives
$\mathcal{L}_X = d\iota_X + \iota_X d$; with $d\omega = 0$ this yields
$\mathcal{L}_X\omega = d\iota_X\omega$, so $\iota_X\omega$ and $\iota_Y\omega$
are closed. Using
$\iota_{[X,Y]} = \mathcal{L}_X\iota_Y - \iota_Y\mathcal{L}_X$ on $\omega$,
\[
\iota_{[X,Y]}\omega
= \mathcal{L}_X(\iota_Y\omega) - \iota_Y(\mathcal{L}_X\omega)
= \bigl(d\iota_X + \iota_X d\bigr)(\iota_Y\omega) - 0
= d\bigl(\iota_X\iota_Y\omega\bigr)
= d\bigl(\omega(Y,X)\bigr),
\]
where $d(\iota_Y\omega)=0$ was used.
(iv) Every $X_f$ satisfies $\mathcal{L}_{X_f}\omega = d\,df = 0$, so (iii)
applies and gives
$\iota_{[X_f,X_g]}\omega = d(\omega(X_g,X_f)) = d\{g,f\} = -d\{f,g\}
= -\iota_{X_{\{f,g\}}}\omega$; nondegeneracy gives the identity.
(v) Using (i) and (iv),
\[
\{f,\{g,h\}\} = X_{\{g,h\}}(f) = -[X_g,X_h](f)
 = -X_g\bigl(X_h(f)\bigr) + X_h\bigl(X_g(f)\bigr)
 = -\{\{f,h\},g\} + \{\{f,g\},h\}.
\]
Antisymmetry, applied twice on the left factor and once on the right, gives
$\{\{f,h\},g\} = -\{g,\{f,h\}\} = \{g,\{h,f\}\}$ and
$\{\{f,g\},h\} = -\{h,\{f,g\}\}$. Substituting,
$\{f,\{g,h\}\} = -\{g,\{h,f\}\} - \{h,\{f,g\}\}$, which is (v).
\end{proof}

\begin{remark}[Sign conventions are part of the type]
\label{rem:signs}
Item (iv) carries a minus sign under the conventions
$\iota_{X_H}\omega = dH$ and $\{f,g\}=\omega(X_f,X_g)$ fixed above, so
$f \mapsto X_f$ is a Lie algebra \emph{anti}-homomorphism and $f\mapsto -X_f$
is a homomorphism. A check on $(\Rr^2,dq\wedge dp)$ with $f=\tfrac12 q^2$,
$g=\tfrac12 p^2$ gives $X_f = -q\partial_p$, $X_g = p\partial_q$,
$[X_f,X_g] = -q\partial_q + p\partial_p$, and $\{f,g\}=qp$ with
$X_{qp} = q\partial_q - p\partial_p$. A compositional reading that treats
$f\mapsto X_f$ as a functor must carry this sign in the data.
\end{remark}

\begin{proposition}[Coordinate form; conservation; Liouville]
\label{prop:coords}
Let $(q^1,\dots,q^n,p_1,\dots,p_n)$ be local coordinates in which
$\omega = \sum_i dq^i\wedge dp_i$. Then
\[
  X_H = \sum_{i=1}^{n}\Bigl(\frac{\partial H}{\partial p_i}\,
  \frac{\partial}{\partial q^i}
  -\frac{\partial H}{\partial q^i}\,\frac{\partial}{\partial p_i}\Bigr),
  \qquad
  \dot q^{\,i} = \frac{\partial H}{\partial p_i},\quad
  \dot p_i = -\frac{\partial H}{\partial q^i},
\]
$\{f,g\} = \sum_i\bigl(\partial_{q^i}f\,\partial_{p_i}g -
\partial_{p_i}f\,\partial_{q^i}g\bigr)$, and moreover
$X_H(H)=0$ and $\mathcal{L}_{X_H}\omega = 0$, hence
$\mathcal{L}_{X_H}\omega^{\wedge n} = 0$.
\end{proposition}

\begin{proof}
Write $X = \sum_i (A^i\partial_{q^i} + B_i\partial_{p_i})$. Then
$\iota_X\omega = \sum_i (A^i\,dp_i - B_i\,dq^i)$. Equating with
$dH = \sum_i(\partial_{q^i}H\,dq^i + \partial_{p_i}H\,dp_i)$ gives
$A^i = \partial_{p_i}H$ and $B_i = -\partial_{q^i}H$, which is the displayed
field and, applied to a curve, Hamilton's equations. The bracket formula is
Lemma~\ref{lem:poisson}(i) in these coordinates.
$X_H(H) = dH(X_H) = \omega(X_H,X_H) = 0$ by antisymmetry.
$\mathcal{L}_{X_H}\omega = d\iota_{X_H}\omega + \iota_{X_H}d\omega
= d\,dH + 0 = 0$, and the Lie derivative is a derivation of the wedge
product, so the volume form $\omega^{\wedge n}$ is preserved.
\end{proof}

\begin{established}[Darboux]
\label{est:darboux}
Under Axioms~\ref{ax:kin}--\ref{ax:symp}, every point of $M$ has a
neighbourhood carrying coordinates in which
$\omega = \sum_{i=1}^{n} dq^i\wedge dp_i$
\textup{\cite[Ch.~3]{AbrahamMarsden1978}}, \textup{\cite[\S43]{Arnold1989}}.
\end{established}

\begin{theorem}[Local reconstruction]
\label{thm:reconstruction}
Let $\Th$ satisfy Axioms~\ref{ax:kin}--\ref{ax:law}. Then every point of $M$
has a neighbourhood $V$ and coordinates on $V$ in which the observable type is
$\Cinf{V,\Rr}$ with the canonical Poisson bracket, the dynamical law is
Hamilton's equations for $H|_V$, and the flow preserves the canonical volume.
Conversely, any system of first-order ODEs on an open $V\subseteq\Rr^{2n}$
that takes the form $\dot q = \partial_p H$, $\dot p = -\partial_q H$ for some
$H \in \Cinf{V,\Rr}$ arises from the object
$(V, \sum_i dq^i\wedge dp_i, \Cinf{V,\Rr}, H, \Phi)$ satisfying
Axioms~\ref{ax:kin}--\ref{ax:law}.
\end{theorem}

\begin{proof}
The forward direction combines Established result~\ref{est:darboux} with
Proposition~\ref{prop:coords}. For the converse, $\omega_0 =
\sum_i dq^i\wedge dp_i$ is closed (constant coefficients) and nondegenerate
(its matrix in the frame $\partial_{q},\partial_p$ is the standard symplectic
matrix, which is invertible), so Axioms~\ref{ax:kin}--\ref{ax:symp} hold;
Axiom~\ref{ax:obs} is a definition; $H$ supplies Axiom~\ref{ax:gen}; and the
computation in Proposition~\ref{prop:coords}, read backwards, identifies the
given ODE with the flow equation of $X_H$, giving Axiom~\ref{ax:law}.
\end{proof}

\begin{remark}[What Theorem~\ref{thm:reconstruction} does \emph{not} say]
\label{rem:notsay}
It is a statement about local normal form. It does not assert that a physical
system is described by such an object, and it supplies no criterion for
deciding whether a given empirical dynamics admits a Hamiltonian
presentation. That decision problem is the inverse problem of the calculus of
variations, solved for two-dimensional second-order systems by Douglas
\cite{Douglas1941} and open in general.
\end{remark}

\subsection{How much of \texorpdfstring{$H$}{H} the dynamics determines}

\begin{theorem}[Exact sequence of the kinematic type]
\label{thm:exact}
Let $\Xf_{\Symp}(M) := \{X : \mathcal{L}_X\omega = 0\}$ and
$\Xf_{\Ham}(M) := \{X_f : f\in\mathcal{A}(M)\}$. Then
\[
0 \longrightarrow H^{0}_{\dR}(M)
  \longrightarrow \mathcal{A}(M)
  \stackrel{f\mapsto X_f}{\longrightarrow} \Xf_{\Symp}(M)
  \stackrel{X\mapsto[\iota_X\omega]}{\longrightarrow} H^{1}_{\dR}(M)
  \longrightarrow 0
\]
is an exact sequence of real vector spaces. Moreover $\Xf_{\Ham}(M)$ is a Lie
ideal in $\Xf_{\Symp}(M)$ and
$[\Xf_{\Symp}(M),\Xf_{\Symp}(M)] \subseteq \Xf_{\Ham}(M)$.
\end{theorem}

\begin{proof}
$H^0_{\dR}(M)$ is the space of locally constant functions. If $X_f = 0$ then
$df = \iota_{X_f}\omega = 0$, so $f$ is locally constant; conversely a locally
constant $f$ has $df=0$, hence $X_f=0$ by Lemma~\ref{lem:XH}. This gives
exactness at $H^0_{\dR}(M)$ and at $\mathcal{A}(M)$.
Every $X_f$ satisfies $\mathcal{L}_{X_f}\omega = d\,df = 0$, so the second map
lands in $\Xf_{\Symp}(M)$. For $X \in \Xf_{\Symp}(M)$, Cartan's formula and
$d\omega=0$ give $d(\iota_X\omega) = \mathcal{L}_X\omega = 0$, so
$[\iota_X\omega]$ is defined. The class vanishes exactly when
$\iota_X\omega = df$ for some $f$, that is exactly when $X = X_f$; this is
exactness at $\Xf_{\Symp}(M)$. For surjectivity onto $H^1_{\dR}(M)$: given a
closed $1$-form $\alpha$, put $X := (\omega^{\flat})^{-1}\alpha$; then
$\iota_X\omega=\alpha$ is closed, so $\mathcal{L}_X\omega = d\alpha = 0$ and
$X\in\Xf_{\Symp}(M)$ with $[\iota_X\omega]=[\alpha]$.
The last claim is Lemma~\ref{lem:poisson}(iii): for
$X,Y \in \Xf_{\Symp}(M)$ the form $\iota_{[X,Y]}\omega = d(\omega(Y,X))$ is
exact, so $[X,Y] = X_{g}$ with $g := \omega(Y,X)$ by Lemma~\ref{lem:XH}; in
particular $[X,Y]\in\Xf_{\Ham}(M)$, and taking $Y \in \Xf_{\Ham}(M)$ shows
$\Xf_{\Ham}(M)$ is an ideal.
\end{proof}

\begin{corollary}[Underdetermination of the generator]
\label{cor:underdet}
If $M$ is connected, the flow $\Phi$ determines $H$ up to one additive real
constant. If $H^{1}_{\dR}(M)\neq 0$, there exist $\omega$-preserving flows
that are generated by no global $H$ whatsoever.
\end{corollary}

\begin{proof}
The first claim is exactness at $\mathcal{A}(M)$ together with
$H^0_{\dR}(M)\cong\Rr$ for connected $M$: two generators of the same field
differ by a locally constant, hence constant, function. The second is
exactness at $\Xf_{\Symp}(M)$: any $X$ with $[\iota_X\omega]\neq 0$ preserves
$\omega$ and is not of the form $X_f$.
\end{proof}

\begin{nonexample}[Symplectic but not Hamiltonian]
\label{nex:torus}
Let $M = \mathbb{T}^2 = \Rr^2/(2\pi\mathbb{Z})^2$ with angle coordinates
$(\theta^1,\theta^2)$ and $\omega = d\theta^1\wedge d\theta^2$, which is
closed and nondegenerate. Let $X = \partial_{\theta^1}$. Then
$\iota_X\omega = d\theta^2$, which is closed but not exact, since
$\int_{\gamma} d\theta^2 = 2\pi \neq 0$ for $\gamma$ the second generating
circle whereas the integral of an exact form over a loop vanishes. Hence $X$
preserves $\omega$ and the Liouville area, its flow is a one-parameter group
of symplectomorphisms, and by Corollary~\ref{cor:underdet} no
$H\in\Cinf{\mathbb{T}^2}$ has $X = X_H$. Axiom~\ref{ax:law} therefore excludes
a dynamics that satisfies every conservation property usually cited as the
physical content of Hamiltonian mechanics.
\end{nonexample}

\begin{interpretation}[Reading of Corollary~\ref{cor:underdet}]
\label{int:zeroenergy}
The additive constant in $H$ is invisible to $\mathrm{Emp}(\Th)$ whenever the
readout $\mu$ is built from the flow, because $H$ and $H+c$ generate the same
$X_H$ by Lemma~\ref{lem:XH}. Within the scope of this paper, an absolute zero
of energy is therefore not an empirically realized element of the theory
object. That classical gravitation and general relativity change this
situation is noted as a boundary fact and is out of scope.
\end{interpretation}

\subsection{Symmetry}

\begin{proposition}[Hamiltonian reciprocity, and conservation]
\label{prop:noether}
For $F,H\in\mathcal{A}(M)$ the following are equivalent:
\textup{(a)} $\{F,H\}=0$;
\textup{(b)} $F$ is constant along the flow of $X_H$;
\textup{(c)} $H$ is constant along the flow of $X_F$.
\end{proposition}

\begin{proof}
By Lemma~\ref{lem:poisson}(i), $\frac{d}{dt}\bigl(F\circ\Phi^{H}_t\bigr)
= X_H(F)\circ\Phi^H_t = \{F,H\}\circ\Phi^H_t$, giving (a)$\Leftrightarrow$(b)
because a smooth function on a connected time interval is constant exactly
when its derivative vanishes identically. Likewise
$\frac{d}{dt}(H\circ\Phi^F_t) = \{H,F\}\circ\Phi^F_t$, and
$\{H,F\} = -\{F,H\}$, giving (a)$\Leftrightarrow$(c).
\end{proof}

The group-theoretic strengthening --- for a Lie group $G$ acting by
symplectomorphisms with an $\mathrm{Ad}^{*}$-equivariant momentum map
$J:M\to\mathfrak{g}^{*}$ and $G$-invariant $H$, each component of $J$ is
conserved --- is an established result we use only by citation
\cite[Chs.~11--13]{MarsdenRatiu1999}, \cite[Ch.~4]{AbrahamMarsden1978}.

% ===============================================================
\section{Composition, and what composition cannot do}
\label{sec:composition}
% ===============================================================

\subsection{Morphisms}

\begin{definition}[Morphism, isomorphism]
\label{def:morphism}
A \emph{morphism} $\varphi : (M,\omega,H) \to (M',\omega',H')$ of bare objects
is a smooth map $\varphi: M\to M'$ with
$\varphi^{*}\omega' = \omega$ and $H = H'\circ\varphi + c$ for some constant
$c\in\Rr$. It is an \emph{isomorphism} when $\varphi$ is in addition a
diffeomorphism. An \emph{isomorphism of theory objects} is an isomorphism of
bare objects together with a bijection of interpretations
$\mathcal{I}\to\mathcal{I}'$ intertwining $P,\mu,\mathcal{U}$.
\end{definition}

\begin{lemma}[Symplectic maps are immersions]
\label{lem:immersion}
If $\varphi^{*}\omega' = \omega$ with $\omega$ nondegenerate, then
$d\varphi_x$ is injective for every $x\in M$.
\end{lemma}

\begin{proof}
Suppose $d\varphi_x(v) = 0$. For every $w \in T_xM$,
$\omega_x(v,w) = (\varphi^{*}\omega')_x(v,w)
= \omega'_{\varphi(x)}\bigl(d\varphi_x v,\,d\varphi_x w\bigr) = 0$.
Nondegeneracy of $\omega_x$ gives $v = 0$.
\end{proof}

\begin{newresult}[No symplectic coarse-graining]
\label{new:coarse}
There is no morphism $\varphi:(M,\omega,H)\to(M',\omega',H')$ in the sense of
Definition~\ref{def:morphism} with $\dim M > \dim M'$. In particular,
projecting a many-degree-of-freedom object onto a few-degree-of-freedom
object is never a morphism of Hamiltonian theory objects.
\end{newresult}

\begin{proof}
By Lemma~\ref{lem:immersion}, $d\varphi_x$ is injective, so
$\dim M = \dim T_xM \le \dim T_{\varphi(x)}M' = \dim M'$.
\end{proof}

\begin{remark}
\label{rem:coarsefolklore}
Lemma~\ref{lem:immersion} is standard and New result~\ref{new:coarse} is
elementary, so we expect the latter is folklore; the claim here is to its
\emph{explicit statement in the typed setting}, because it settles a question
a compositional account must answer. Coarse-graining, thermodynamic reduction
and effective-model extraction are \emph{not} morphisms in this category. They
are quotient constructions --- symplectic reduction
\cite{MarsdenRatiu1999} --- so a compositional framework that wants
them must carry a second class of arrows, or work in a double category or a
category of spans.
\end{remark}

\begin{lemma}[Isomorphisms intertwine flows]
\label{lem:intertwine}
If $\varphi$ is an isomorphism then $d\varphi\circ X_H = X_{H'}\circ\varphi$
and $\varphi\circ\Phi_t = \Phi'_t\circ\varphi$ on the common domain of
definition.
\end{lemma}

\begin{proof}
For $w\in T_xM$,
$\omega'_{\varphi(x)}(d\varphi\, X_H,\ d\varphi\, w)
= (\varphi^{*}\omega')_x(X_H,w) = \omega_x(X_H,w) = dH_x(w)
= d(H'\circ\varphi)_x(w) = dH'_{\varphi(x)}(d\varphi\,w)$.
As $\varphi$ is a diffeomorphism, $d\varphi_x$ is onto $T_{\varphi(x)}M'$, so
$\iota_{d\varphi X_H}\omega' = dH'$ at $\varphi(x)$, and
Lemma~\ref{lem:XH} gives $d\varphi\,X_H = X_{H'}\circ\varphi$. Two vector
fields that are $\varphi$-related have $\varphi$-related flows.
\end{proof}

\subsection{The monoidal product and its exact range}

\begin{definition}[Cartesian product of objects]
\label{def:product}
For bare objects $\Th_i^{\mathrm{math}} = (M_i,\omega_i,\mathcal{A}(M_i),H_i,\Phi^i)$,
$i=1,2$, set
\[
  \Th_1^{\mathrm{math}}\otimes\Th_2^{\mathrm{math}}
  := \bigl(M_1\times M_2,\ \pi_1^{*}\omega_1+\pi_2^{*}\omega_2,\
  \mathcal{A}(M_1\times M_2),\ H_1\circ\pi_1 + H_2\circ\pi_2,\ \Phi\bigr).
\]
\end{definition}

\begin{proposition}[$\otimes$ is well typed and non-interacting]
\label{prop:productwelltyped}
$\omega := \pi_1^{*}\omega_1+\pi_2^{*}\omega_2$ satisfies
Axiom~\ref{ax:symp} on $M_1\times M_2$, and with
$H_\oplus := H_1\circ\pi_1 + H_2\circ\pi_2$ one has
$X_{H_\oplus} = X_{H_1}\oplus X_{H_2}$ and
$\Phi_t = \Phi^1_t\times\Phi^2_t$.
\end{proposition}

\begin{proof}
$d\omega = \pi_1^{*}d\omega_1 + \pi_2^{*}d\omega_2 = 0$.
Under the canonical splitting
$T_{(x_1,x_2)}(M_1\times M_2)\cong T_{x_1}M_1\oplus T_{x_2}M_2$ the form
$\omega$ is the block-diagonal pairing $\omega_1\oplus\omega_2$, and a
block-diagonal bilinear form with nondegenerate blocks is nondegenerate:
if $\omega((v_1,v_2),(w_1,w_2)) = \omega_1(v_1,w_1)+\omega_2(v_2,w_2)$
vanishes for all $(w_1,w_2)$, then taking $w_2=0$ gives $v_1 = 0$ and taking
$w_1=0$ gives $v_2=0$. For the vector field, with
$X := X_{H_1}\oplus X_{H_2}$,
$\iota_X\omega = \pi_1^{*}(\iota_{X_{H_1}}\omega_1)
+ \pi_2^{*}(\iota_{X_{H_2}}\omega_2)
= \pi_1^{*}dH_1 + \pi_2^{*}dH_2 = dH_\oplus$,
and Lemma~\ref{lem:XH} identifies $X = X_{H_\oplus}$. The flow of a direct-sum
field on a product is the product of the flows.
\end{proof}

\begin{newresult}[Interaction is outside the range of $\otimes$]
\label{new:interaction}
Let $M = M_1\times M_2$ carry $\omega$ and $H_\oplus$ as above and let
$H = H_\oplus + V$ with $V \in \mathcal{A}(M)$. Then the following are
equivalent:
\begin{enumerate}[label=(\alph*),leftmargin=*,itemsep=1pt]
\item $\Phi^{H}_t = \Phi^{1}_t\times\Phi^{2}_t$ for all $t$ in a neighbourhood
      of $0$ and all points of $M$;
\item $X_V = 0$;
\item $V$ is locally constant.
\end{enumerate}
Consequently the image of $\otimes$ inside the class of objects on
$M_1\times M_2$ with the split symplectic form consists exactly of the
composites whose factor dynamics are independent; no choice of factors makes
$\otimes$ produce an interacting composite.
\end{newresult}

\begin{proof}
(c)$\Rightarrow$(b): $dV=0$ gives $X_V = (\omega^{\flat})^{-1}(0) = 0$.
(b)$\Rightarrow$(a): $X_H = X_{H_\oplus} + X_V$ by linearity of
$f\mapsto X_f$ (Lemma~\ref{lem:XH} and linearity of $d$ and of
$(\omega^{\flat})^{-1}$); with $X_V=0$,
Proposition~\ref{prop:productwelltyped} gives the product flow.
(a)$\Rightarrow$(b): differentiating $\Phi^H_t(x)$ and
$(\Phi^1_t\times\Phi^2_t)(x)$ at $t=0$ gives
$X_H(x) = (X_{H_1}\oplus X_{H_2})(x) = X_{H_\oplus}(x)$ for every $x$, hence
$X_V = X_H - X_{H_\oplus} = 0$.
(b)$\Rightarrow$(c): $X_V = 0$ gives $\iota_{X_V}\omega = dV = 0$ by
Lemma~\ref{lem:XH}, so $V$ is locally constant.
For the final sentence: an object on $(M_1\times M_2,\omega)$ with generator
$H$ lies in the image of $\otimes$ exactly when $H - H_\oplus$ is locally
constant for some $H_1,H_2$, which by the equivalence is exactly when its flow
splits.
\end{proof}

\begin{corollary}[A compositional account needs a second operation]
\label{cor:secondop}
Under Definition~\ref{def:product}, the assignment
$(\Th_1,\Th_2)\mapsto\Th_1\otimes\Th_2$ is a symmetric monoidal product on
bare objects with unit the one-point object $(\{*\},0,\Rr,0)$, and by
New result~\ref{new:interaction} every interacting system on
$M_1\times M_2$ is outside its range. Hence no account of physical
composition that uses only $\otimes$ can express interaction.
\end{corollary}

\begin{proof}
Symmetry and associativity up to the canonical diffeomorphisms of products are
inherited from the Cartesian product of manifolds; each canonical
diffeomorphism pulls back the split form to the split form and matches the
generators, so it is an isomorphism in the sense of
Definition~\ref{def:morphism}. Degeneracy of $\omega$ is not an issue for the
unit because $\dim\{*\} = 0$ and the empty pairing is vacuously
nondegenerate. The second sentence is
New result~\ref{new:interaction}. Combining gives the conclusion.
\end{proof}

\begin{remark}[Relation to the applied-category-theory literature]
\label{rem:act}
Corollary~\ref{cor:secondop} is the precise form of a widely acknowledged
point: interconnection is modelled by cospans, decorated cospans, or
relations, not by monoidal product alone
\cite{Fong2015,BaezStay2011,MacLane1998}. Our contribution here is only the
sharp ``if and only if'' of New result~\ref{new:interaction} inside the typed
Hamiltonian setting, which converts a methodological preference into a
theorem about what $\otimes$ can represent.
\end{remark}

\subsection{The natural repair fails to be a category}

\begin{established}[The symplectic ``category'' is not a category]
\label{est:weinstein}
Take symplectic manifolds as objects and Lagrangian submanifolds of
$M_1^{-}\times M_2$ (canonical relations) as arrows. The set-theoretic
composite of two canonical relations need not be a submanifold, and when it
is, need not be Lagrangian; composability requires transversality and
clean-intersection hypotheses. Consequently this data does not form a
category without further work; see Weinstein \cite{Weinstein2010} for the
statement of the difficulty and for the standard repairs.
\end{established}

\begin{interpretation}[Cost of Established result~\ref{est:weinstein}]
\label{int:weinstein}
The natural candidate for a composition rich enough to carry interaction ---
composition of canonical relations --- is only \emph{partially} defined. A
compositional foundation for classical mechanics must therefore accept a
partial composition, pass to a completion, or restrict to a subclass on which
transversality is guaranteed. This is a structural limit on the thesis under
test, not a gap in exposition.
\end{interpretation}

% ===============================================================
\section{Empirical realization, equivalence, and limits}
\label{sec:empirical}
% ===============================================================

\subsection{Two directions of a single separation}

\begin{definition}[$\varepsilon$-empirical equivalence]
\label{def:eeq}
Let $\Th,\Th'$ be theory objects with the same $M$, the same preparable set
$P$, the same readout $\mu$ and units, resolution $\varepsilon$, and window
$[0,T]$. They are \emph{$\varepsilon$-empirically equivalent on $(P,T)$} when
\[
  \sup_{z\in P}\ \sup_{t\in[0,T]}\
  \bigl\|\mu(\Phi_t z) - \mu(\Phi'_t z)\bigr\| \;\le\; \varepsilon .
\]
For objects on different manifolds the comparison requires an explicit
translation of preparations and readouts; we do not use that case.
\end{definition}

\begin{proposition}[Empirically equivalent, mathematically non-isomorphic]
\label{prop:eqnotiso}
Fix $m,k>0$, $R>0$, $T>0$, and on $(\Rr^2,\omega_0=dq\wedge dp)$ set
\[
  H_0 = \frac{p^2}{2m}+\frac{k}{2}q^2, \qquad
  H_\lambda = H_0 + \lambda q^4 \quad (0<\lambda\le 1),
  \qquad P = \{z : \|z\|\le R\}.
\]
Then:
\begin{enumerate}[label=(\roman*),leftmargin=*,itemsep=1pt]
\item there is $C = C(m,k,R,T)<\infty$, independent of $\lambda$, with
      $\sup_{z\in P}\sup_{t\in[0,T]}\|\Phi^{\lambda}_t z - \Phi^{0}_t z\|
      \le C\lambda$; hence for every $\varepsilon>0$ and every $1$-Lipschitz
      readout $\mu$ there is $\lambda>0$ making $\Th_0$ and $\Th_\lambda$
      $\varepsilon$-empirically equivalent on $(P,T)$;
\item for every $\lambda>0$ the bare objects $(\Rr^2,\omega_0,H_0)$ and
      $(\Rr^2,\omega_0,H_\lambda)$ are not isomorphic in the sense of
      Definition~\ref{def:morphism}.
\end{enumerate}
\end{proposition}

\begin{proof}
(i) Put $E_R := \sup_{\|z\|\le R} H_1(z) < \infty$. For $\lambda\in(0,1]$ and
$z\in P$ one has $H_\lambda(z)\le H_1(z)\le E_R$, and $H_\lambda$ is conserved
along $\Phi^\lambda$ by Proposition~\ref{prop:coords}. Since
$H_0 \le H_\lambda$ pointwise, the trajectory satisfies
$H_0(\Phi^\lambda_t z) \le E_R$, so it remains in the compact set
$K := \{H_0\le E_R\}$; the same bound holds for $\lambda = 0$. Set
$q_{\max} := \sqrt{2E_R/k}$ and $A_R := 4q_{\max}^3$, an upper bound on
$\|X_{q^4}\|$ on $K$ (indeed $X_{q^4} = -4q^3\partial_p$). The field
$X_{H_0}(q,p) = (p/m,\,-kq)$ is linear, hence globally Lipschitz with some
constant $L>0$. Writing $u(t) := \Phi^{\lambda}_t z - \Phi^{0}_t z$,
\[
  \dot u = \bigl[X_{H_0}(\Phi^\lambda_t z)-X_{H_0}(\Phi^0_t z)\bigr]
          + \lambda X_{q^4}(\Phi^\lambda_t z),
  \qquad \|\dot u\| \le L\|u\| + \lambda A_R,\quad u(0)=0 .
\]
Gr\"onwall's inequality gives
$\|u(t)\| \le (\lambda A_R/L)(e^{Lt}-1)$, so
$C := (A_R/L)(e^{LT}-1)$ works. A $1$-Lipschitz $\mu$ transports the bound,
and $\lambda < \varepsilon/C$ gives the equivalence.

(ii) For $\lambda \ge 0$ the potential $V_\lambda(q)=\tfrac k2 q^2+\lambda q^4$
is smooth, even, strictly increasing on $[0,\infty)$ with
$V_\lambda(q)\to\infty$, and $V'_\lambda(q)= kq+4\lambda q^3$ vanishes only at
$q=0$. Hence for every $E>0$ the level set $\{H_\lambda = E\}$ is a regular,
compact, connected curve, and the orbit on it is periodic with period
$T_\lambda(E)$. Let $A_\lambda(E)$ be the area enclosed. The classical
relation $T_\lambda(E) = A_\lambda'(E)$ for one-degree-of-freedom systems with
regular compact level sets \cite[Ch.~10]{Arnold1989} is used.

For $\lambda = 0$ the level set is the ellipse with semi-axes
$\sqrt{2E/k}$ and $\sqrt{2mE}$, so $A_0(E) = 2\pi E\sqrt{m/k}$ and
$T_0(E) = 2\pi\sqrt{m/k}$ for every $E>0$: the period function is constant.

For $\lambda>0$ suppose, for contradiction, that
$T_\lambda \equiv c$ on $(0,\infty)$ for some $c>0$. Since $A_\lambda$ is
continuous with $A_\lambda(0)=0$ and $A_\lambda' = T_\lambda$, integration
gives $A_\lambda(E) = cE$ for all $E \ge 0$. On the other hand
$H_\lambda \ge \frac{p^2}{2m}+\lambda q^4$ pointwise, so
$\{H_\lambda\le E\}\subseteq\{\frac{p^2}{2m}+\lambda q^4\le E\}$ and
\[
  A_\lambda(E) \;\le\;
  \int_{-\sqrt{2mE}}^{\sqrt{2mE}}
  2\Bigl(\frac{E-\frac{p^2}{2m}}{\lambda}\Bigr)^{1/4} dp
  \;=\; 2\lambda^{-1/4}\sqrt{2mE}\;E^{1/4}\!\!\int_{-1}^{1}\!(1-u^2)^{1/4}du
  \;=\; \kappa\,\sqrt{m}\,\lambda^{-1/4}E^{3/4},
\]
with $\kappa := 2\sqrt{2}\int_{-1}^{1}(1-u^2)^{1/4}du \in (0,\infty)$, using
the substitution $p=\sqrt{2mE}\,u$. Combining,
$cE \le \kappa\sqrt m\,\lambda^{-1/4}E^{3/4}$, i.e.
$E^{1/4}\le \kappa\sqrt m/(c\lambda^{1/4})$ for all $E>0$, which fails for
$E$ large. Hence $T_\lambda$ is not constant.

Now suppose $\varphi$ were an isomorphism $(\Rr^2,\omega_0,H_0)\to
(\Rr^2,\omega_0,H_\lambda)$. By Lemma~\ref{lem:intertwine} it carries orbits
of $\Phi^0$ to orbits of $\Phi^{\lambda}$ preserving the time parameter, so
it carries periodic orbits to periodic orbits of the same period, and it is
onto. Every periodic orbit of $\Phi^{\lambda}$ with energy $E>0$ is therefore
the image of a periodic orbit of $\Phi^0$, whose period is
$2\pi\sqrt{m/k}$. That makes $T_\lambda$ constant on $(0,\infty)$, contrary to
the previous paragraph.
\end{proof}

\begin{proposition}[Mathematically isomorphic, physically distinct]
\label{prop:isonotphys}
There exist theory objects $\Th,\Th'$ with $U\Th \cong U\Th'$ as bare objects
whose empirical models differ, and no isomorphism of bare objects induces a
map $\mathrm{Emp}(\Th)\to\mathrm{Emp}(\Th')$.
\end{proposition}

\begin{proof}
Take $\Th$ the mass--spring object on $(\Rr^2,dq\wedge dp)$ with
$H = p^2/2m + kq^2/2$, readout $\mu(q,p)=q$ with $\mathcal{U}(\mu)$ the metre
and $\mathcal{U}(H)$ the joule, and $P$ the states reachable by displacing the
mass. Take $\Th'$ the ideal $LC$ circuit on $(\Rr^2,dQ\wedge d\Phi_{\!f})$
with $H' = \Phi_{\!f}^2/2L + Q^2/2C$, readout $\mu'(Q,\Phi_{\!f})=Q$ with
$\mathcal{U}(\mu')$ the coulomb, and $P'$ the states reachable by charging the
capacitor. Choosing $m = L$ and $k = 1/C$, the identity map of $\Rr^2$ is an
isomorphism of bare objects: it pulls back the form to the form and matches
the generators exactly.

The empirical models differ as sets-with-units:
$\mathrm{Emp}(\Th)$ is a set of position records in metres and
$\mathrm{Emp}(\Th')$ a set of charge records in coulombs, and
$\mathcal{U}$ is part of the data of an interpretation by
Definition~\ref{def:interpretation}. A bare isomorphism is a map of manifolds
and carries no component acting on $\mathcal{U}$, because $U$ deletes
$\mathcal{U}$; hence it induces no map of empirical models. Two further
differences are recorded as empirical inputs rather than derived here: the
admissible interventions differ (one may clamp the spring; one may open the
circuit), and the failure modes differ (plastic deformation; dielectric
breakdown).
\end{proof}

\begin{interpretation}[The separation, stated]
\label{int:separation}
Propositions~\ref{prop:eqnotiso} and \ref{prop:isonotphys} together show that
mathematical isomorphism of bare Hamiltonian objects is neither sufficient
(\ref{prop:isonotphys}) nor necessary (\ref{prop:eqnotiso}) for empirical
indistinguishability. Any claim of the form ``these two theories are the same
physics because the mathematics is the same'' therefore requires a separate
empirical argument about $P$, $\mu$, $\mathcal{U}$, $\varepsilon$ and $T$. We
adopt this as a working discipline rather than as a discovery: the pair of
counterexamples is what makes the discipline enforceable.
\end{interpretation}

\subsection{Limiting relations}

\begin{example}[Small-oscillation limit is not an isomorphism]
\label{ex:pendulum}
Let $\Th^{\mathrm{pend}}$ be the planar pendulum on
$T^{*}S^{1}$ with $H^{\mathrm{pend}}(\theta,p_\theta)=
\frac{p_\theta^2}{2m\ell^2}+mg\ell(1-\cos\theta)$, and let
$\Th^{\mathrm{ho}}$ be the harmonic object on $T^{*}\Rr$ with
$H^{\mathrm{ho}}=\frac{p^2}{2m\ell^2}+\frac{1}{2}mg\ell\,\theta^2$. Since
$1-\cos\theta = \tfrac12\theta^2 + O(\theta^4)$, on
$\{|\theta|\le\delta,\ |p_\theta|\le\delta\}$ the two generators differ by
$O(\delta^4)$ and a Gr\"onwall estimate of the form used in
Proposition~\ref{prop:eqnotiso}(i) bounds the trajectory difference on
$[0,T]$ by $C(T)\delta^4$. The two objects are not isomorphic: the pendulum
has a non-constant period function and, unlike the harmonic object, has
unbounded ($\theta$-circulating) orbits and a separatrix. The relation
between them is an \emph{asymptotic} one, indexed by $\delta$, not an
isomorphism of theory objects.
\end{example}

\begin{interpretation}[Reading of Example~\ref{ex:pendulum}]
\label{int:limits}
Limiting relations between theory objects are relations between
\emph{interpreted} objects at stated resolution and time window; they are not
recovered by any operation on bare objects alone, because the statement
``$C(T)\delta^4 \le \varepsilon$'' quantifies over $\varepsilon$ and $T$,
which live in $\mathcal{I}$.
\end{interpretation}

\subsection{An empirical anchor and an empirical refutation}

\begin{empirical}[Anchor]
\label{emp:anchor}
For laboratory mass--spring and $LC$ systems in regimes where damping,
nonlinearity and drift are below the readout resolution, the object of
Proposition~\ref{prop:isonotphys} reproduces the recorded time series within
$\varepsilon$ over $[0,T]$. This is taken as an empirical input to the pilot;
no new measurement is reported in this paper.
\end{empirical}

\begin{established}[Refutation of a specific Hamiltonian object]
\label{est:mercury}
The two-body Newtonian gravitational object on $T^{*}(\Rr^3\setminus\{0\})$
with $H = \frac{\|\mathbf{p}\|^2}{2\mu}-\frac{GM\mu}{\|\mathbf{q}\|}$
predicts a closed, non-precessing ellipse. The observed perihelion advance of
Mercury exceeds the Newtonian $N$-body accounting by an amount matched by the
general-relativistic prediction of $42.98''$ per century
\cite{Will2014}. The Hamiltonian object above is therefore empirically false
at that precision.
\end{established}

\begin{interpretation}[Reading of Established result~\ref{est:mercury}]
\label{int:mercury}
Established result~\ref{est:mercury} refutes one \emph{object}, not the
\emph{type}: the first post-Newtonian correction is itself standardly
presented by a Hamiltonian on a finite-dimensional phase space, so the failure
was of the generator $H$ and not of Axioms~\ref{ax:kin}--\ref{ax:law}. This
distinction --- refuting a member versus refuting a type --- is what the typed
presentation is for; Section~\ref{sec:falsification} states what would refute
the type.
\end{interpretation}

% ===============================================================
\section{Non-examples: where each axiom bites}
\label{sec:nonexamples}
% ===============================================================

\begin{nonexample}[Odd dimension: Axiom~\ref{ax:kin} bites]
\label{nex:odd}
No manifold of odd dimension admits a nondegenerate $2$-form. If
$\dim V = d$ is odd and $\beta$ is an alternating bilinear form on $V$, its
matrix $B$ in any basis satisfies $B^{\!\top} = -B$, so
$\det B = \det(-B^{\!\top}) = (-1)^{d}\det B = -\det B$, giving $\det B = 0$
and hence a nonzero $v$ with $\beta(v,\cdot)=0$. Systems whose natural state
space is odd-dimensional --- for instance a state space carrying a contact
structure, or a thermodynamic state space with an extra scalar --- fall
outside Axiom~\ref{ax:kin} and require a different kinematic type.
\end{nonexample}

\begin{nonexample}[Degenerate bracket: Axiom~\ref{ax:symp} bites]
\label{nex:rigidbody}
On $\mathfrak{so}(3)^{*}\cong\Rr^3$ the Lie--Poisson bracket
$\{f,g\}(\mathbf{x}) = -\,\mathbf{x}\cdot(\nabla f\times\nabla g)$ is a
Poisson bracket for which the free rigid body,
$H(\mathbf{x}) = \tfrac12\sum_i x_i^2/I_i$, generates Euler's equations
\cite{MarsdenRatiu1999}. The associated bivector has rank $2$
at $\mathbf{x}\neq 0$ and rank $0$ at the origin, so it is not induced by any
nondegenerate $2$-form; the function $\|\mathbf{x}\|^2$ is a Casimir, i.e.
$\{\|\mathbf{x}\|^2,\cdot\} = 0$, which is impossible for a nonconstant
function under Axiom~\ref{ax:symp} because $\{f,\cdot\}=0$ with
$\{f,g\} = X_g(f) = -X_f(g)$ for all $g$ forces $X_f = 0$, hence $df=0$. The
object is Poisson but not symplectic; the theory type must be widened, and
widening it changes which statements of Section~\ref{sec:reconstruction}
survive.
\end{nonexample}

\begin{newresult}[Dissipation: Axiom~\ref{ax:law} bites, for every $\omega$]
\label{new:damped}
Let $\gamma>0$, $\omega_0^2>0$ and let $X$ be the damped-oscillator field on
$\Rr^2$,
\[
  X(x,v) \;=\; \bigl(v,\ -\omega_0^2 x-\gamma v\bigr).
\]
Then there is \emph{no} symplectic form $\sigma$ on $\Rr^2$ and no
$H\in\Cinf{\Rr^2}$ with $\iota_X\sigma = dH$. In particular the obstruction is
not an artefact of the standard form $dx\wedge dv$.
\end{newresult}

\begin{proof}
The linear part of $X$ has matrix
$\left(\begin{smallmatrix} 0 & 1\\ -\omega_0^2 & -\gamma\end{smallmatrix}\right)$
with characteristic polynomial $s^2+\gamma s+\omega_0^2$, whose roots have
real part $-\gamma/2 < 0$ (for real roots, both are negative because their sum
is $-\gamma<0$ and product $\omega_0^2>0$). As $X$ is linear, every solution
is a combination of $e^{s_\pm t}$ modes, so $\Phi_t(z)\to 0$ as $t\to+\infty$
for every $z\in\Rr^2$.

Suppose $\iota_X\sigma = dH$ for some symplectic $\sigma$ and some smooth $H$.
Then $X(H) = dH(X) = \sigma(X,X) = 0$ by antisymmetry, so $H$ is constant
along every trajectory: $H(\Phi_t z) = H(z)$ for all $t\ge 0$. Letting
$t\to+\infty$ and using continuity of $H$ gives $H(z) = H(0)$ for every
$z\in\Rr^2$. Hence $dH \equiv 0$, so $\iota_X\sigma \equiv 0$, and
nondegeneracy of $\sigma$ gives $X\equiv 0$. This contradicts
$X(1,0) = (0,-\omega_0^2)\neq 0$.
\end{proof}

\begin{remark}[Precision of New result~\ref{new:damped}]
\label{rem:damped}
Three qualifications keep the statement honest.
(a) The proof uses global asymptotic stability, hence $\gamma>0$,
$\omega_0^2>0$ and the global domain $\Rr^2$; on a proper subdomain avoiding
the attractor, or with a \emph{time-dependent} structure, Hamiltonian
presentations of damped systems exist and are standard.
(b) Being Hamiltonian is thus a property of the pair $(X,\sigma)$ together
with the domain, not of $X$ alone; which second-order systems admit a
variational or Hamiltonian presentation is the inverse problem of the calculus
of variations \cite{Douglas1941}.
(c) The result is elementary and the argument is likely known; the claim is to
the explicit form-independent statement, which the typed setting requires and
which is stronger than the usual divergence argument --- the latter excludes
only the standard $\omega_0$, since $\dive X = -\gamma \neq 0$.
\end{remark}

\begin{established}[No functorial quantization: the type does not extend]
\label{est:gvh}
There is no map from the Poisson algebra of polynomials on $\Rr^{2n}$ to
self-adjoint operators on a Hilbert space that is linear, sends $1$ to the
identity, sends $\{f,g\}$ to $\frac{1}{i\hbar}[\hat f,\hat g]$, and acts
irreducibly on the linear polynomials. This is the Groenewold--van Hove
obstruction; statements, proofs, and the analogous results for $S^2$ and
$T^{*}S^1$ are collected by Gotay \cite{Gotay1998}.
\end{established}

\begin{interpretation}[Reading of Established result~\ref{est:gvh}]
\label{int:gvh}
A compositional programme that seeks a functor from classical theory objects
to quantum theory objects preserving the full observable algebra structure
cannot succeed on $\Rr^{2n}$: the obstruction is a theorem, not a missing
construction. Quantum mechanics is a boundary test for this paper, and we
claim nothing inside it; we record the obstruction only because it bounds
what the reconstruction can be extended to.
\end{interpretation}

\begin{established}[Symplectic rigidity]
\label{est:gromov}
If $r>R$ there is no symplectic embedding of the ball
$B^{2n}(r)$ into the cylinder $B^{2}(R)\times\Rr^{2n-2}$, although a
volume-preserving embedding exists for all $r,R>0$ when $n\ge 2$
\cite{Gromov1985}.
\end{established}

\begin{interpretation}[Reading of Established result~\ref{est:gromov}]
\label{int:gromov}
Axiom~\ref{ax:symp} is strictly stronger than ``the flow preserves phase
volume''. The kinematic type therefore carries information that the measure
type does not, and the choice between them is in principle a substantive
choice rather than a notational one.
\end{interpretation}

\begin{conjecture}[Empirical bite of rigidity]
\label{conj:rigidity}
There exist a finite-dimensional classical laboratory system, an
interpretation $\mathcal{I}=(P,\mu,\mathcal{U},\varepsilon,T)$, and a pair of
candidate objects agreeing on all volume-theoretic predictions within
$\varepsilon$ on $(P,T)$ whose predictions for $\mu$ differ by more than
$\varepsilon$ solely as a consequence of Established
result~\ref{est:gromov}. We know of no such realization and do not assert that
one exists.
\end{conjecture}

\begin{interpretation}[Reading of Established result~\ref{est:darboux}]
\label{int:darboux}
Established result~\ref{est:darboux} entails that symplectic manifolds of
equal dimension have no local invariants: any two are locally isomorphic.
Hence all local empirical content of a Hamiltonian theory object resides in $H$ and
in $\mathcal{I}$, never in $(M,\omega)$; the kinematic type contributes only
global information (topology of $M$, cohomology classes as in
Theorem~\ref{thm:exact} and Non-example~\ref{nex:torus}) and the dimension
count. A compositional account that locates physical content in the kinematic
type must therefore locate it globally, or the content is empty.
\end{interpretation}

% ===============================================================
\section{Falsification conditions and status of the thesis}
\label{sec:falsification}
% ===============================================================

\begin{definition}[What would refute the type]
\label{def:refute}
We say the Hamiltonian type is \emph{refuted for a domain} $D$ of physical
systems when there is a system in $D$, an interpretation $\mathcal{I}$, and
recorded data $\mathcal{D}$ such that no object satisfying
Axioms~\ref{ax:kin}--\ref{ax:law} on any finite-dimensional $M$ has
$\mathcal{D}\subseteq_{\varepsilon}\mathrm{Emp}(\Th)$.
\end{definition}

Three concrete refuting observations, each of which is currently unavailable
and each of which is stated as a target rather than a claim:

\begin{enumerate}[label=(F\arabic*),leftmargin=*,itemsep=2pt]
\item \textbf{Bracket failure.} Measured response functions for a closed
      finite-dimensional system whose induced bracket on the measured
      observables violates the Jacobi identity beyond $\varepsilon$. By
      Lemma~\ref{lem:poisson}(v) this refutes $d\omega=0$, hence
      Axiom~\ref{ax:symp}, for that system.
\item \textbf{Global generator failure.} A closed finite-dimensional system
      whose measured evolution preserves $\omega$ but for which the class
      $[\iota_X\omega]\in H^1_{\dR}(M)$ is measurably nonzero. By
      Corollary~\ref{cor:underdet} and Non-example~\ref{nex:torus} this
      refutes Axiom~\ref{ax:law} while leaving Axioms~\ref{ax:kin}--\ref{ax:obs}
      intact --- a discriminating rather than a global refutation.
\item \textbf{Irreducible dissipation.} A system that is closed under the
      operative notion of isolation, is finite-dimensional, and exhibits
      global asymptotic stability of an equilibrium. By
      New result~\ref{new:damped} this refutes Axiom~\ref{ax:law} for that
      system under \emph{every} symplectic form on its state space.
\end{enumerate}

\begin{interpretation}[Status of the pilot thesis after this paper]
\label{int:status}
The reconstruction supports the following bounded report, and nothing
stronger.
\emph{Precision:} the thesis can be stated precisely for this domain, by
Definition~\ref{def:object}, Axioms~\ref{ax:kin}--\ref{ax:interp} and
Definition~\ref{def:interpretation}.
\emph{Internal coherence:} the object is consistent and generates the standard
formalism (Theorems~\ref{thm:reconstruction} and \ref{thm:exact}).
\emph{Compatibility:} it reproduces the established structural results cited
here without modification.
\emph{Usefulness:} the typing did work --- it produced
New results~\ref{new:coarse}, \ref{new:interaction} and \ref{new:damped}, and
forced the two-directional separation of
Interpretation~\ref{int:separation}.
\emph{Costs:} the compositional reading is incomplete in a specific, provable
way. $\otimes$ cannot express interaction
(New result~\ref{new:interaction}); the natural richer composition is only
partial (Established result~\ref{est:weinstein}); the category has no
coarse-graining arrows (New result~\ref{new:coarse}); and the type does not
extend functorially to quantum theory (Established result~\ref{est:gvh}).
We therefore do not report the thesis as established. We report that, for
finite-dimensional Hamiltonian mechanics, it is precise, non-trivial, and
partially falsifiable, with four named structural deficits.
\end{interpretation}

\begin{conjecture}[Interconnection completeness]
\label{conj:interconnection}
There is a symmetric monoidal double category whose loose arrows are
partially-defined canonical relations and whose tight arrows are
symplectomorphisms, in which every object
$(M_1\times M_2,\ \omega_1\oplus\omega_2,\ H_1\oplus H_2 + V)$ arises as a
composite of $\Th_1,\Th_2$ and one arrow determined by $V$. Decorated cospans
\cite{Fong2015} and the monoidal apparatus of
\cite{MacLane1998,BaezStay2011} are the natural tools; the obstruction to a
naive construction is Established result~\ref{est:weinstein}.
\end{conjecture}

% ===============================================================
\appendix

\section{Type skeleton in Haskell}
\label{app:haskell}

The following is a rendering of Definitions~\ref{def:types},
\ref{def:object}, \ref{def:interpretation} and \ref{def:morphism} as
types. It is a \emph{typing}, not a formal proof: it makes the arities and
dependencies of the theory object machine-checkable, and it makes
New result~\ref{new:interaction} expressible as a property, but nothing here
verifies any theorem of this paper. No claim of formal verification is made.

\begin{lstlisting}[language=Haskell]
module Principia.Hamiltonian where
type Tangent = [Double];  type Bilinear = (Tangent, Tangent) -> Double
type Observable m = m -> Double            -- observable type
type Flow m = Double -> m -> Maybe m       -- Maybe encodes the partiality

-- Kinematic type (kinematic + symplectic axioms); 'm' = points of M.
data Kinematic m = Kinematic { dimHalf :: Int, symplForm :: m -> Bilinear }
newtype Generator m = Generator (Observable m)

-- Interpretation. Units are DATA, not derivable from the bare object;
-- that is exactly what the "isomorphic but distinct" result turns on.
data Unit = Metre | Coulomb | Joule | Second | Other String deriving (Eq, Show)
data Interp m = Interp { preparable :: m -> Bool, readout :: m -> [Double]
                       , units :: [Unit], resolution, window :: Double }

data BareObject m   = BareObject (Kinematic m) (Generator m) (Flow m)
data TheoryObject m = TheoryObject (BareObject m) (Interp m)
forget :: TheoryObject m -> BareObject m               -- the operation U
forget (TheoryObject b _) = b

-- Cartesian monoidal product: block-diagonal form, summed generator,
-- paired flow.  'split' evaluates omega_1 (+) omega_2 in the splitting.
tensor :: BareObject a -> BareObject b -> BareObject (a, b)
tensor (BareObject ka (Generator ha) fa) (BareObject kb (Generator hb) fb) =
  BareObject (Kinematic (dimHalf ka + dimHalf kb) split)
             (Generator (\(x,y) -> ha x + hb y))
             (\t (x,y) -> (,) <$> fa t x <*> fb t y)
  where split (x,y) (v,w) = let n = 2 * dimHalf ka
                                (v1,v2) = splitAt n v; (w1,w2) = splitAt n w
                            in symplForm ka x (v1,w1) + symplForm kb y (v2,w2)

-- The non-interaction result as a sampled PREDICATE, not a theorem: a
-- coupling is non-interacting exactly when its differential vanishes.
-- This names the criterion; it does not prove the equivalence.
newtype Coupling m = Coupling (Observable m)
isNonInteracting :: (Observable m -> m -> [Double]) -> Coupling m -> [m] -> Bool
isNonInteracting grad (Coupling v) zs = all (all ((< 1e-12) . abs) . grad v) zs

-- Morphisms: the type records the data, it does not enforce the
-- immersion lemma, which is proved in the text and not here.
data Morphism a b = Morphism { mapPoints :: a -> b, shiftConst :: Double }
\end{lstlisting}

\section{Source verification record}
\label{app:verification}

\footnotesize
Author list, title, year, and venue were matched against publisher or
repository records for every cited item before citation. Confirmed without
material residual uncertainty: \cite{BaezStay2011} (LNP 813, 2011,
pp.~95--174), \cite{Douglas1941} (Trans.\ AMS \textbf{50}, 1941, 71--128),
\cite{Fong2015} (TAC \textbf{30}, 2015, no.~33, 1096--1120),
\cite{Gromov1985} (Invent.\ Math.\ \textbf{82}, 1985, 307--347, DOI
10.1007/BF01388806), \cite{MacLane1998} (Springer GTM 5, 2nd ed., 1998),
\cite{Will2014} (Living Rev.\ Relativ.\ \textbf{17}, 2014, 4;
arXiv:1403.7377, where the figure $42.98''$/century appears).
Confirmed with a stated residual: \cite{AbrahamMarsden1978}
(Benjamin/Cummings, 2nd ed., 1978, ISBN 978-0-8053-0102-1 --- chapter
pointers are to that edition and no page numbers are asserted);
\cite{Arnold1989} (Springer GTM 60, 2nd ed., 1989, trans.\ Vogtmann \&
Weinstein --- the pointers $\S43$ and Ch.~10 were not independently
re-paginated); \cite{MarsdenRatiu1999} (Springer TAM 17, 2nd ed., 1999, ISBN
978-0-387-98643-2 --- chapter pointers to that edition);
\cite{Gotay1998} (arXiv:math-ph/9809011, 1998 --- the preprint is cited
because the published-volume pagination was not confirmed);
\cite{Weinstein2010} (Port.\ Math.\ \textbf{67}, 2010, start page 261
confirmed --- the end page 278 is inferred and not confirmed).

\footnotesize
\begin{thebibliography}{99}
\setlength{\itemsep}{0pt}

\bibitem{AbrahamMarsden1978}
R.~Abraham and J.~E. Marsden,
\emph{Foundations of Mechanics}, 2nd ed., revised, enlarged and reset, with
the assistance of T.~Ratiu and R.~Cushman,
Benjamin/Cummings, Reading, MA, 1978.

\bibitem{Arnold1989}
V.~I. Arnold,
\emph{Mathematical Methods of Classical Mechanics}, 2nd ed.,
translated by K.~Vogtmann and A.~Weinstein,
Graduate Texts in Mathematics 60, Springer, New York, 1989.

\bibitem{BaezStay2011}
J.~C. Baez and M.~Stay,
\emph{Physics, topology, logic and computation: a Rosetta Stone}, in
B.~Coecke (ed.), New Structures for Physics,
Lecture Notes in Physics 813, Springer, Berlin, 2011, pp.~95--174;
arXiv:0903.0340.

\bibitem{Douglas1941}
J.~Douglas,
\emph{Solution of the inverse problem of the calculus of variations},
Trans.\ Amer.\ Math.\ Soc.\ \textbf{50} (1941), 71--128.

\bibitem{Fong2015}
B.~Fong,
\emph{Decorated cospans},
Theory Appl.\ Categ.\ \textbf{30} (2015), no.~33, 1096--1120;
arXiv:1502.00872.

\bibitem{Gotay1998}
M.~J. Gotay,
\emph{Obstructions to quantization},
arXiv:math-ph/9809011 (1998).

\bibitem{Gromov1985}
M.~Gromov,
\emph{Pseudo holomorphic curves in symplectic manifolds},
Invent.\ Math.\ \textbf{82} (1985), 307--347.

\bibitem{MacLane1998}
S.~Mac~Lane,
\emph{Categories for the Working Mathematician}, 2nd ed.,
Graduate Texts in Mathematics 5, Springer, New York, 1998.

\bibitem{MarsdenRatiu1999}
J.~E. Marsden and T.~S. Ratiu,
\emph{Introduction to Mechanics and Symmetry: A Basic Exposition of Classical
Mechanical Systems}, 2nd ed.,
Texts in Applied Mathematics 17, Springer, New York, 1999.

\bibitem{Weinstein2010}
A.~Weinstein,
\emph{Symplectic categories},
Port.\ Math.\ \textbf{67} (2010), 261--278.

\bibitem{Will2014}
C.~M. Will,
\emph{The confrontation between general relativity and experiment},
Living Rev.\ Relativ.\ \textbf{17} (2014), 4; arXiv:1403.7377.

\end{thebibliography}

\end{document}
